\section{矩阵运算的扰动分析}
\label{sec:multivariable-matrix-operation-perturbation}

第~\ref{chap:matrix-differentiation}章已经求出常用矩阵运算的微分，并用增量的一阶项验证了可微性。本节回到误差问题：矩阵数据带有扰动$\bm E$时，这些微分给出的线性预测在有限步长下有多准，误差在最坏方向上被放大多少。

对每种运算都要同时检查三件事：输出是矩阵还是标量，因而误差用什么范数或尺度度量；扰动是否受对称等结构限制；整条扰动线段是否留在定义域内。前一节的一般原则在这里落到具体的余项表达式上。本节的矩阵范数除特别说明外取满足乘积不等式$\lVert\bm X\bm Y\rVert\leq\lVert\bm X\rVert\lVert\bm Y\rVert$的范数，例如谱范数或Frobenius范数；求逆时取由向量范数诱导的矩阵范数，它还满足$\lVert\bm I\rVert=1$。

\subsection{基本矩阵运算的增量与余项}
\label{sec:multivariable-basic-operation-remainder}

第~\ref{sec:multivariable-matrix-basic-operations}节对转置、乘法、平方、迹、Frobenius范数平方和二次型写出了精确增量。这些增量都是扰动的多项式，所以余项不必另行估计，直接从增量中读出，再用范数控制大小。

\paragraph{线性运算}
加法、数乘、转置、迹以及$\bm X\mapsto\bm A\bm X\bm B$都是线性映射，精确增量就是一阶项，余项恒为零。对这些运算，线性预测在任意步长下都精确，误差问题只剩放大倍数。由定理~\ref{thm:multivariable-local-condition-number}，局部绝对条件数就是运算本身的算子范数，且与基点无关。

例如在Frobenius范数下，转置不改变任何矩阵的长度，条件数为$1$。迹则不同：由Cauchy--Schwarz不等式，
\[
 |\operatorname{tr}\bm E|=|\langle\bm I,\bm E\rangle_F|\leq\sqrt d\,\lVert\bm E\rVert_F ,
\]
取$\bm E=\bm I/\sqrt d$时等号成立，所以$d\times d$矩阵的迹的条件数为$\sqrt d$。各对角元的误差同号时会累加起来。

\paragraph{乘积与平方}
由式~\eqref{eq:multivariable-product-increment}，两个因子同时扰动时，
\[
 (\bm A+\bm E)(\bm B+\bm H)-\bm A\bm B-(\bm A\bm H+\bm E\bm B)=\bm E\bm H .
\]
由乘积不等式，一阶项与余项分别满足
\[
 \begin{gathered}
 \lVert\bm A\bm H+\bm E\bm B\rVert\leq\lVert\bm A\rVert\lVert\bm H\rVert+\lVert\bm E\rVert\lVert\bm B\rVert,\\
 \lVert\bm E\bm H\rVert\leq\lVert\bm E\rVert\lVert\bm H\rVert .
 \end{gathered}
\]
两者的比较说明：两个扰动都不超过$\rho$时，一阶项至多与$\rho$同阶，余项至多是$\rho^2$。

平方映射同理，$(\bm X+\bm E)^2-\bm X^2$的一阶项为$\bm X\bm E+\bm E\bm X$，余项恰为$\bm E^2$，且$\lVert\bm E^2\rVert\leq\lVert\bm E\rVert^2$。

这些界可以直接用作误差预算。若要求线性预测的绝对误差不超过$\tau$，乘积情形只需$\lVert\bm E\rVert\lVert\bm H\rVert\leq\tau$，两个扰动同样大小时步长不超过$\sqrt\tau$即可。若要求余项不超过一阶项的某个比例，则还要知道一阶项在所取方向上有多大，不能只看扰动的大小。
\begin{example}[矩阵平方的一阶预测]
\label{ex:multivariable-square-first-order-check}
取$\bm X=\operatorname{diag}(2,1)$与扰动$\bm E=\varepsilon\bm K$，其中
\[
 \bm K=\begin{bmatrix}0&1\\1&0\end{bmatrix}.
\]
一阶项与余项分别为
\[
 \bm X\bm E+\bm E\bm X=\varepsilon\begin{bmatrix}0&3\\3&0\end{bmatrix},\qquad
 \bm E^2=\varepsilon^2\bm I ,
\]
谱范数分别为$3\varepsilon$与$\varepsilon^2$。取$\varepsilon=0.1$，一阶项为$0.3$，余项为$0.01$，约为一阶项的$3\%$。

若要求余项不超过一阶项的$1\%$，需要
\[
 \frac{\varepsilon^2}{3\varepsilon}\leq0.01 ,
\]
即扰动须满足$\varepsilon\leq0.03$。
\end{example}

\paragraph{Frobenius范数平方与二次型}
标量输出的运算同样可以读出余项。Frobenius范数平方的精确增量为
\[
 \lVert\bm X+\bm E\rVert_F^2-\lVert\bm X\rVert_F^2=2\langle\bm X,\bm E\rangle_F+\lVert\bm E\rVert_F^2 ,
\]
余项恰为$\lVert\bm E\rVert_F^2$，与方向无关。一阶项$2\langle\bm X,\bm E\rangle_F$却可能为零：扰动与$\bm X$正交时，线性预测说“没有变化”，实际变化却是$\lVert\bm E\rVert_F^2$。这时余项与一阶项之比没有意义，应直接给出绝对误差的预算。

二次型$\bm x^\top\bm A\bm x$在$\bm A$与$\bm x$同时扰动时，第~\ref{sec:multivariable-matrix-basic-operations}节写出的余项为
\[
 \bm h^\top\bm A\bm h+\bm x^\top\bm E\bm h+\bm h^\top\bm E\bm x+\bm h^\top\bm E\bm h .
\]
在谱范数下，它的绝对值不超过
\[
 \lVert\bm A\rVert_2\lVert\bm h\rVert_2^2+2\lVert\bm E\rVert_2\lVert\bm x\rVert_2\lVert\bm h\rVert_2+\lVert\bm E\rVert_2\lVert\bm h\rVert_2^2 .
\]
每一项都含两个扰动因子，其中$\lVert\bm E\rVert_2\lVert\bm h\rVert_2$一项还乘了基点的大小$\lVert\bm x\rVert_2$：基点越大，矩阵与向量扰动的交叉作用越强。

比较起来，线性运算的余项为零，多项式运算的余项是扰动的二次或更高次项，并且可以精确写出。误差预算由这些高次项决定线性预测可用的步长；导数本身的大小只决定一阶项，不决定预测还差多少。

\subsection{求逆的扰动}
\label{sec:multivariable-inverse-perturbation}

第~\ref{sec:multivariable-inverse-differential}节由逆差恒等式证明了求逆可微，并得到$\mathrm d(\bm X^{-1})=-\bm X^{-1}\mathrm d\bm X\bm X^{-1}$。那里只用到余项是$o(\lVert\bm E\rVert)$。本节把余项精确写出，给出它的二次界，再计算求逆的条件数，并说明它与线性方程解映射的联系。

\paragraph{精确余项与可逆邻域}
设$\bm A$可逆，$\bm B=\bm A+\bm E$，$q=\lVert\bm A^{-1}\rVert\lVert\bm E\rVert<1$。由逆差恒等式~\eqref{eq:ch06-inverse-finite-perturbation}，
\[
 \bm B^{-1}-\bm A^{-1}=-\bm B^{-1}\bm E\bm A^{-1}.
\]
把右端的$\bm B^{-1}$写成$\bm A^{-1}+(\bm B^{-1}-\bm A^{-1})$，再用一次同一恒等式，
\[
 \bm B^{-1}-\bm A^{-1}
 =-\bm A^{-1}\bm E\bm A^{-1}-(\bm B^{-1}-\bm A^{-1})\bm E\bm A^{-1}
 =-\bm A^{-1}\bm E\bm A^{-1}+\bm B^{-1}\bm E\bm A^{-1}\bm E\bm A^{-1}.
\]
第一项是一阶预测，第二项是精确余项。由$\lVert\bm B^{-1}\rVert\leq\lVert\bm A^{-1}\rVert/(1-q)$（第~\ref{sec:ls-inverse-sensitivity}节），
\begin{equation}
 \lVert\bm B^{-1}\bm E\bm A^{-1}\bm E\bm A^{-1}\rVert
 \leq\lVert\bm B^{-1}\rVert\lVert\bm A^{-1}\rVert^2\lVert\bm E\rVert^2
 \leq\frac{\lVert\bm A^{-1}\rVert^3\lVert\bm E\rVert^2}{1-q}.
\label{eq:multivariable-inverse-remainder-bound}
\end{equation}

第~\ref{chap:linear-recurrence-stability}章的Neumann级数给出同一结论的另一种看法。写$\bm B=\bm A(\bm I-\bm T)$，其中$\bm T=-\bm A^{-1}\bm E$，$\lVert\bm T\rVert\leq q<1$，因而谱半径小于$1$。由定理~\ref{thm:ls-neumann}，$\bm I-\bm T$可逆，且
\[
 \bm B^{-1}=\sum_{k\geq0}\bm T^k\bm A^{-1}.
\]
前两项$\bm A^{-1}-\bm A^{-1}\bm E\bm A^{-1}$正是一阶预测；截断误差界~\eqref{eq:ls-neumann-tail}取$N=1$，给出余项不超过$q^2\lVert\bm A^{-1}\rVert/(1-q)$，与式~\eqref{eq:multivariable-inverse-remainder-bound}相同。级数还说明，余项依次由$\bm E$的二次、三次……项组成。

\paragraph{一阶预测的误差}
式~\eqref{eq:multivariable-inverse-remainder-bound}同时是有限扰动下一阶预测的误差界。对$0\leq t\leq1$，$\lVert\bm A^{-1}(t\bm E)\rVert\leq tq<1$，所以整条扰动线段$\bm A+t\bm E$都由可逆矩阵组成，定义域条件在整段上成立。

于是可以比较三个量：实际逆矩阵$(\bm A+\bm E)^{-1}$，一阶预测$\bm A^{-1}-\bm A^{-1}\bm E\bm A^{-1}$，以及二者之差。第~\ref{sec:ls-inverse-sensitivity}节控制的是全变化，
\[
 \lVert(\bm A+\bm E)^{-1}-\bm A^{-1}\rVert\leq\frac{\lVert\bm A^{-1}\rVert^2\lVert\bm E\rVert}{1-q},
\]
与$\lVert\bm E\rVert$同阶；本段新增的是从全变化中扣除一阶项之后，近似还差多少，它与$\lVert\bm E\rVert^2$同阶。$q$接近$1$时两种界都变得很大，一阶预测不再可靠。
\begin{example}[逆矩阵的一阶预测与余项]
\label{ex:multivariable-inverse-first-order-check}
取$\bm A=\operatorname{diag}(2,1)$与扰动
\[
 \bm E=\begin{bmatrix}0&0.1\\0.1&0\end{bmatrix},
\]
采用谱范数，则$\lVert\bm A^{-1}\rVert_2=1$、$\lVert\bm E\rVert_2=0.1$、$q=0.1$。这个矩阵可以直接求逆，
\[
 (\bm A+\bm E)^{-1}=\frac1{1.99}\begin{bmatrix}1&-0.1\\-0.1&2\end{bmatrix},
\]
而一阶预测为
\[
 \bm A^{-1}-\bm A^{-1}\bm E\bm A^{-1}=\begin{bmatrix}0.5&-0.05\\-0.05&1\end{bmatrix}.
\]
实际值减去预测，剩下的余项为
\[
 \frac{0.0005}{1.99}\begin{bmatrix}10&-1\\-1&20\end{bmatrix}\approx\begin{bmatrix}0.00251&-0.00025\\-0.00025&0.00503\end{bmatrix}.
\]

现在比较几个量的大小。余项的谱范数约为$0.00505$，约为一阶项谱范数$0.05$的十分之一，并落在式~\eqref{eq:multivariable-inverse-remainder-bound}的上界$0.01/0.9\approx0.0111$之内。

再看全变化。全变化$(\bm A+\bm E)^{-1}-\bm A^{-1}$的谱范数约为$0.0540$，第~\ref{sec:ls-inverse-sensitivity}节的全变化界为$0.1/0.9\approx0.111$。
\end{example}
一阶预测扣除了全变化中约$0.05$的主要部分，剩下的误差约为扰动大小的平方量级；余项界比实际误差大约两倍，它只保证误差不超过多少。

\paragraph{求逆的条件数}
在无结构限制的扰动模型下，输入、输出都取谱范数。由乘积不等式，
\[
 \lVert\bm A^{-1}\bm E\bm A^{-1}\rVert_2\leq\lVert\bm A^{-1}\rVert_2^2\lVert\bm E\rVert_2 .
\]

这个上界可以取到。设$\sigma_d$是$\bm A$的最小奇异值，对应的左、右奇异向量为$\bm u_d,\bm v_d$，即$\bm A\bm v_d=\sigma_d\bm u_d$、$\bm A^\top\bm u_d=\sigma_d\bm v_d$（第~\ref{sec:ls-svd}章）。取$\bm E=\bm u_d\bm v_d^\top$，它的谱范数为$1$。由$\bm A^{-1}\bm u_d=\bm v_d/\sigma_d$与$\bm A^{-\top}\bm v_d=\bm u_d/\sigma_d$，
\[
 \bm A^{-1}\bm E\bm A^{-1}=(\bm A^{-1}\bm u_d)(\bm A^{-\top}\bm v_d)^\top=\frac{\bm v_d\bm u_d^\top}{\sigma_d^2},
\]
其谱范数为$1/\sigma_d^2=\lVert\bm A^{-1}\rVert_2^2$。

由定理~\ref{thm:multivariable-local-condition-number}，求逆映射在$\bm A$处的局部绝对条件数为$\lVert\bm A^{-1}\rVert_2^2$，相对条件数为
\[
 \lVert\bm A^{-1}\rVert_2^2\cdot\frac{\lVert\bm A\rVert_2}{\lVert\bm A^{-1}\rVert_2}=\lVert\bm A\rVert_2\lVert\bm A^{-1}\rVert_2=\kappa_2(\bm A).
\]

矩阵条件数在这里重新出现，但含义不同：这里输入是矩阵本身，输出是逆矩阵，而第~\ref{sec:ls-matrix-condition-number}节的输入是向量。若限制扰动结构，例如只允许对称扰动，须改用微分在允许子空间上的范数。求解方程的算法及其数值稳定性仍按第~\ref{chap:solution-error-regularization}章处理。

\paragraph{线性方程的解映射}
求逆与乘法复合，得到线性方程的解映射
\[
 \operatorname{sol}(\bm A,\bm b)=\bm A^{-1}\bm b ,
\]
它定义在可逆$\bm A$与任意$\bm b$组成的开集上。固定可逆的$\bm A$，记$\bm x=\bm A^{-1}\bm b$；联合输入$(\bm E,\bm f)$取乘积范数$\lVert(\bm E,\bm f)\rVert=(\lVert\bm E\rVert^2+\lVert\bm f\rVert^2)^{1/2}$，其中$\bm f$是右端的增量。

第~\ref{sec:ls-equation-sensitivity}节从扰动后的方程
\[
 (\bm A+\bm E)(\bm x+\Delta\bm x)=\bm b+\bm f
\]
中减去$\bm A\bm x=\bm b$，得到
\[
 (\bm A+\bm E)\Delta\bm x=\bm f-\bm E\bm x .
\]
在$q<1$时，这给出精确式
\[
 \Delta\bm x=(\bm A+\bm E)^{-1}(\bm f-\bm E\bm x).
\]
把$(\bm A+\bm E)^{-1}$拆成$\bm A^{-1}$与逆差，
\[
 \Delta\bm x=\bm A^{-1}(\bm f-\bm E\bm x)+\bigl((\bm A+\bm E)^{-1}-\bm A^{-1}\bigr)(\bm f-\bm E\bm x).
\]
第一项对$(\bm E,\bm f)$线性。第二项中，两个因子分别满足
\[
 \begin{gathered}
 \bigl\lVert(\bm A+\bm E)^{-1}-\bm A^{-1}\bigr\rVert\leq\frac{\lVert\bm A^{-1}\rVert^2\lVert\bm E\rVert}{1-q},\\
 \lVert\bm f-\bm E\bm x\rVert\leq\lVert\bm f\rVert+\lVert\bm E\rVert\lVert\bm x\rVert ;
 \end{gathered}
\]
二者都不超过联合扰动大小的常数倍，所以第二项是联合扰动的二次量。因此解映射可微，且
\begin{equation}
 \mathrm d\operatorname{sol}_{(\bm A,\bm b)}(\bm E,\bm f)=\bm A^{-1}(\bm f-\bm E\bm x).
\label{eq:multivariable-solution-map-differential}
\end{equation}
也可以用链式法则核对：$\operatorname{sol}$是求逆与矩阵—向量乘法的复合，由乘积的联合微分与式~\eqref{eq:analysis-inverse-differential}，一阶项为
\[
 (-\bm A^{-1}\bm E\bm A^{-1})\bm b+\bm A^{-1}\bm f=\bm A^{-1}(\bm f-\bm E\bm x),
\]
结果相同。第~\ref{sec:multivariable-implicit-function}节已用隐函数公式核对过同一结果。

第~\ref{sec:ls-equation-sensitivity}节的界控制完整的解误差，式~\eqref{eq:multivariable-solution-map-differential}则提取局部的一阶预测，上面的拆分还说明预测的误差是联合扰动的二次量。数据误差、残差与算法稳定性的完整分析仍以第~\ref{chap:solution-error-regularization}章为准，本节只补充一阶响应及其精度。

\subsection{行列式的扰动}
\label{sec:multivariable-determinant-perturbation}

第~\ref{sec:multivariable-determinant-differential}节逐列展开$\det(\bm A+\bm E)$，得到一阶项$\langle\bm C,\bm E\rangle_F$，其中$\bm C$是代数余子式矩阵；展开中其余各项至少含两列扰动。这些高阶项就是线性预测的误差。它们对包括奇异矩阵在内的全部方阵都有意义，所以先讨论绝对误差，再在可逆时讨论相对误差。

\paragraph{绝对变化的余项}
在固定的有界邻域内，余项有二次控制。记$K=\{\bm A+\bm F:\lVert\bm F\rVert_F\leq1\}$。行列式的二阶偏导数都是多项式，在闭有界集$K$上有界，所以存在$M$使$\lVert D^2\det_{\bm Z}\rVert\leq M$对$\bm Z\in K$成立。

对$\lVert\bm E\rVert_F\leq1$，令$g(t)=\det(\bm A+t\bm E)$，它是$t$的多项式。它的二阶导数及其界为
\[
 g''(t)=D^2\det_{\bm A+t\bm E}[\bm E,\bm E],\qquad |g''(t)|\leq M\lVert\bm E\rVert_F^2 .
\]
一元Taylor公式给出
\[
 \bigl|\det(\bm A+\bm E)-\det\bm A-\langle\bm C,\bm E\rangle_F\bigr|\leq\frac M2\lVert\bm E\rVert_F^2 ,
\]
常数$M$依赖所取的邻域。

二阶情形可以精确写出余项。对$\bm A,\bm E\in\mathbb R^{2\times2}$，直接展开得
\[
 \det(\bm A+\bm E)=\det\bm A+\langle\bm C,\bm E\rangle_F+\det\bm E ,
\]
余项就是$\det\bm E=E_{1,1}E_{2,2}-E_{1,2}E_{2,1}$。由$|ab|\leq(a^2+b^2)/2$，
\[
 |\det\bm E|\leq\frac{\lVert\bm E\rVert_F^2}2 ,
\]
即上式中可取$M=1$，而且这个界与基点$\bm A$无关，奇异的$\bm A$也适用。

\paragraph{条件数与相对变化}
对给定扰动方向，一阶响应是$\langle\bm C,\bm E\rangle_F$。由定理~\ref{thm:multivariable-local-condition-number}，局部绝对条件数是这个线性泛函在所选输入范数下的大小，即第~\ref{sec:ls-dual-norms}节的对偶范数：
\[
 \kappa_{\mathrm{abs}}(\det,\bm A)=\max_{\lVert\bm E\rVert\leq1}|\langle\bm C,\bm E\rangle_F| .
\]
输入取Frobenius范数时，由Cauchy--Schwarz不等式，$\bm C\neq\bm0$时最大值在$\bm E=\bm C/\lVert\bm C\rVert_F$处取到，所以$\kappa_{\mathrm{abs}}=\lVert\bm C\rVert_F$，可以直接由代数余子式算出。

可逆时还可以比较相对变化：由$\bm C=\det(\bm A)\bm A^{-\top}$，
\[
 \kappa_{\mathrm{rel}}(\det,\bm A)=\lVert\bm C\rVert_F\frac{\lVert\bm A\rVert_F}{|\det\bm A|}=\lVert\bm A^{-1}\rVert_F\lVert\bm A\rVert_F .
\]
奇异点处$\det\bm A=0$，只有绝对变化有定义，不能除以零行列式。

在相对尺度下，行列式对矩阵变化的反应可以远比矩阵本身的变化剧烈。
\begin{example}[接近奇异时的行列式]
\label{ex:ls-determinant-near-singular}
设$\bm A=\operatorname{diag}(1,\varepsilon)$，$0<\varepsilon<1$，它的行列式$\det\bm A=\varepsilon$很小。只改动第二个对角元，取扰动$\bm E=\operatorname{diag}(0,\delta)$，使它变为$\varepsilon+\delta$。矩阵采用谱范数、标量采用绝对值时，矩阵的相对变化为$\lVert\bm E\rVert_2/\lVert\bm A\rVert_2=|\delta|$；行列式则变为$\varepsilon+\delta$，相对变化为
\[
 \frac{|(\varepsilon+\delta)-\varepsilon|}{\varepsilon}=\frac{|\delta|}{\varepsilon}.
\]
用微分公式计算也得到同样的增量：
\[
 \det(\bm A)\operatorname{tr}(\bm A^{-1}\bm E)=\varepsilon\cdot\frac\delta\varepsilon=\delta .
\]
沿这个方向，行列式是$\delta$的一次函数，所以一阶预测没有余项。

两相比较，行列式的相对变化是矩阵相对变化的$1/\varepsilon$倍。$\varepsilon$越小，矩阵越接近奇异，沿这个方向的相对放大就越强。
\end{example}

但不能仅凭行列式的数值判断求逆是否病态。取$\bm A=\varepsilon\bm I\in\mathbb R^{d\times d}$（$\varepsilon>0$），行列式$\varepsilon^d$可以极小，而由第~\ref{sec:multivariable-inverse-perturbation}节，求逆的相对条件数是$\kappa_2(\varepsilon\bm I)=1$；行列式本身的相对条件数$\lVert\bm A^{-1}\rVert_F\lVert\bm A\rVert_F=d$也不随$\varepsilon$变大。行列式与逆矩阵是两个不同的输出，绝对与相对又是两种尺度，必须分别判断。

\subsection{对数行列式的扰动}
\label{sec:multivariable-logdet-perturbation}

由第~\ref{sec:multivariable-logdet-differential}节，$\log\det$在$\det>0$的开集上可微，微分为$\bm E\mapsto\operatorname{tr}(\bm A^{-1}\bm E)$。它给出一阶变化：
\[
 \log\det(\bm A+\bm E)-\log\det\bm A
 =\operatorname{tr}(\bm A^{-1}\bm E)+o(\lVert\bm E\rVert_F).
\]
相应的局部绝对条件数是线性泛函$\bm E\mapsto\langle\bm A^{-\top},\bm E\rangle_F$在所选输入范数下的对偶范数，Frobenius范数下为$\lVert\bm A^{-1}\rVert_F$。对数行列式可以为零或接近零，这时以它为分母的相对误差没有意义，应直接比较绝对变化。

\paragraph{整条线段的定义域}
要把一阶预测用于有限扰动，先要保证整条线段$\bm A+t\bm E$（$0\leq t\leq1$）留在定义域内。一般的正行列式矩阵，由行列式的连续性，总可以选足够小的邻域，使整条线段$\bm A+t\bm E$仍满足$\det>0$，但这不给出邻域的具体大小。正定情形则可以用第~\ref{chap:condition-and-spectral-perturbation}章的最小特征值界明确检查整条线段保持正定，下面的余项界就基于这一检查。

\paragraph{二次余项界}
设$\bm A$对称正定，最小特征值为$m>0$，$\bm E$对称且$\lVert\bm E\rVert_2<m$。对任意单位向量$\bm x$和$0\leq t\leq1$，
\[
 \bm x^\top(\bm A+t\bm E)\bm x\geq m-\lVert\bm E\rVert_2>0 ,
\]
所以整条线段保持正定，而且最小特征值有统一的正下界。

令$g(t)=\log\det(\bm A+t\bm E)$。由式~\eqref{eq:analysis-logdet-gradient}与链式法则，
\[
 g'(t)=\operatorname{tr}\bigl((\bm A+t\bm E)^{-1}\bm E\bigr);
\]
再用沿参数曲线的求逆公式，
\[
 g''(t)=-\operatorname{tr}\!\left((\bm A+t\bm E)^{-1}\bm E
 (\bm A+t\bm E)^{-1}\bm E\right).
\]
逆矩阵的谱范数至多为$(m-\lVert\bm E\rVert_2)^{-1}$。由Frobenius内积的Cauchy--Schwarz不等式与混合乘积界~\eqref{eq:matrix-analysis-mixed-submultiplicativity}，
\[
 |g''(t)|\leq\lVert(\bm A+t\bm E)^{-1}\bm E\rVert_F^2\leq\frac{\lVert\bm E\rVert_F^2}{(m-\lVert\bm E\rVert_2)^2}.
\]
于是一元Taylor公式给出
\begin{equation}
 \left|\log\det(\bm A+\bm E)-\log\det\bm A
 -\operatorname{tr}(\bm A^{-1}\bm E)\right|
 \leq\frac{\lVert\bm E\rVert_F^2}{2(m-\lVert\bm E\rVert_2)^2}.
\label{eq:multivariable-logdet-remainder-bound}
\end{equation}

例如$\bm A=\operatorname{diag}(2,1)$，$\bm E=\operatorname{diag}(0.1,0)$时，一阶预测为$0.05$，实际变化为$\log(1.05)\approx0.04879$。此时$m=1$，上述上界约为$0.00618$，大于实际误差约$0.00121$。这个界只保证误差不超过多少，不表示误差会接近上界。

导数给出变化的方向与大小，正定线段给出估计的有效范围，余项界说明预测的可靠程度，三者要一起使用。正定条件是这一显式估计的条件，不是一般正行列式矩阵上$\log\det$的定义条件。

上界的分母$m-\lVert\bm E\rVert_2$是扰动后最小特征值的保底下界，即到奇异边界的余量。$\bm A$接近奇异时，$m$很小：一阶响应的条件数$\lVert\bm A^{-1}\rVert_F\geq1/m$变大，余项界又按$1/(m-\lVert\bm E\rVert_2)^2$放大，允许的步长必须小于$m$。

例如$\bm A=\operatorname{diag}(1,\varepsilon)$、$\bm E=\operatorname{diag}(0,\delta)$时，精确变化是$\log(1+\delta/\varepsilon)$，一阶预测是$\delta/\varepsilon$；$\delta$与$\varepsilon$同阶时两者相差显著，$\delta\to-\varepsilon$时精确值趋于$-\infty$。
